\section{Contest setup}

\entry{Header}{paste first}{%
Every entry in this document assumes this block. Ranges are \textbf{half-open}
\code{[l, r)} and graphs are \textbf{0-indexed} unless the entry says otherwise
--- both are stated wherever it matters.}

%HEADER
\begin{lstlisting}
#include <bits/stdc++.h>
#include <cassert>   // bits/stdc++.h no longer pulls this in
using namespace std;
typedef long long ll;
typedef long double ld;
typedef unsigned long long ull;
typedef __int128 i128;
typedef pair<ll, ll> pii;
typedef vector<ll> vi;
typedef vector<vi> matrix;
#define rep(i, a, b) for(ll i = a; i < b; i++)
#define down(i, b, a) for(ll i = b - 1; i >= a; i--)
#define all(v) (v).begin(), (v).end()
const ll INF = 1e18;

int main(){
    ios_base::sync_with_stdio(false);
    cin.tie(NULL);
}
\end{lstlisting}

\entry{Compile and stress test}{}{%
Drop the sanitizers for the real submission --- they cost $\sim$3x.
\code{-Wconversion} additionally catches silent
\code{ll}$\rightarrow$\code{int} narrowing.\\[0.3ex]
Stress test: \code{gen.py} prints a small random case, \code{mine} is the fast
solution, \code{ans} the slow one you trust.}

%NOCOMPILE
\begin{lstlisting}
g++ -std=gnu++17 -O2 -Wall -Wextra -Wshadow \
    -fsanitize=address,undefined -g a.cpp

for i in $(seq 1000); do
  python3 gen.py $i > t.txt
  ./mine < t.txt > o1.txt; ./ans < t.txt > o2.txt
  cmp -s o1.txt o2.txt || { echo "WA $i"; cat t.txt; break; }
done
\end{lstlisting}

\entry{int128 I/O}{}{%
\code{cin} / \code{cout} cannot handle \code{\_\_int128}. Range is
$\pm1.7\cdot10^{38}$, so it is the fix for \code{a*b \% m} when $m$ is up to
$9\cdot10^{18}$: \code{(i128)a * b \% m}.}

\begin{lstlisting}
i128 read128() {
    i128 x = 0; int s = 1; char c = getchar();
    while (c == ' ' || c == '\n') c = getchar();
    if (c == '-') { s = -1; c = getchar(); }
    for (; c >= '0' && c <= '9'; c = getchar())
        x = x * 10 + (c - '0');
    return s * x;
}
void print128(i128 x) {
    if (x < 0) { putchar('-'); x = -x; }
    if (x > 9) print128(x / 10);
    putchar('0' + (int)(x % 10));
}
\end{lstlisting}
